% Ad Quintum Fratrem 2.8 (ad-quintum-fratrem-02-08) --- English translation
% Translated by Claude (Anthropic) from the Latin (Perseus).
% Style guide: STYLE.md.

\ciceroLetterOpener{MARCVS QVINTO FRATRI salutem}

\ciceroSection{1} You are afraid you may interrupt me? In the first place, even if I were absorbed in something, do you, of all people, know what an interruption is — you, or Statius? By Hercules, you seem to be lecturing me on a politeness of just that sort which I, for my part, do not in the least require from you. As for you, I would have you address me and interrupt me and talk back at me and talk with me. What is sweeter to me? By Hercules, no muse-stricken poet \textit{[Greek: mousopataktos]} reads out his latest verses with more relish than I listen to you on any subject at all — public, private, country, town. But it was my own dull modesty that made me fail to take you with me when I left. You set against me, once, an unanswerable \textit{[Greek: anantilekton]} ground — our young Cicero's health; I fell silent.

\ciceroSection{2} A second time, the boys; I gave way. Now this letter of yours, full of charm, has sprinkled into my pleasure this one bit of irritation: that you seem to have been afraid you were a bother to me, and even now to be afraid of it. I should quarrel with you over it, if I might; but by Hercules, if I ever suspect such a thing, I shall say nothing else except that I shall be afraid I may at some time be a bother to you, when I am at your side. I see you have groaned. Yes, that is how it goes — \textit{[Greek corrupt: eid' en aiai ezesas]}; for I shall never say \textit{[Greek corrupt: ea pasas]}. As for our Marius, by Hercules I should have flung him into a litter — not that famous Asician one of King Ptolemy's. For I remember, when it was carrying the man from Naples to Baiae in an eight-bearer Asician litter with a hundred sword-men following, what wonderful laughter we let out when he, all unaware of his escort, suddenly threw the curtains open, and he nearly collapsed from fright and I from laughter. Marius, as I say, I should certainly have brought along, so that I might at last touch some shred of the old urbanity and the most cultivated conversation; but I was unwilling to invite a frail man into a villa still wide open, and not even rough-finished as yet.

\ciceroSection{3} This, however, will be a special pleasure of mine — to enjoy him here as well. For of those estates of mine, you may take it, my neighbour Marius is a beacon. At Anicius's place we shall see to it that everything is made ready; for we are such philologists that we can live even among the carpenters. We get this philosophy not from Hymettus but from \textit{[Greek corrupt: araxira]}. Marius, what is more, is on the weaker side both in health and in temperament.

\ciceroSection{4} As to interruption: I shall take only so much of my time from the two of you for writing as you give me. May you give me none, that I may pause from rest through your fault rather than through any sloth of my own. That you should be too sore at heart over the commonwealth grieves me — that you are a better citizen than Philoctetes, who, after the wrong done him, sought out the very spectacles which I see are bitter to you. I beseech you, fly here (I shall console you and wipe away all your pain) and bring with you, if you love me, Marius; but make haste, both of you. The garden is at home.
